\section{Alternating-group colorings}\label{sec:color}

For any group, the Wirtinger rule at a positive braid crossing is
\(T(a,c)=(c,c^{-1}ac)\); its inverse is
\(T^{-1}(a,c)=(aca^{-1},a)\). A tuple of labels at the top of a closed
braid defines a homomorphism from the knot group exactly when the
crossing rules return the tuple at the bottom. The image is generated
by the top labels, because they generate the Wirtinger group.

Put \(g=\sigma_3\sigma_4\sigma_3^{-1}\).
On the entries \((a,C,e)\) at positions \(3,4,5\), its action is
\[
                 (CeC^{-1},\,C,\,e^{-1}C^{-1}aCe).
\]
In particular, \(g\) fixes the fourth label and acts as \(T\) on
\((a,CeC^{-1})\). The subgroup of \(B_b\) generated by
\(\sigma_2,g,\sigma_5,\ldots,\sigma_{b-1}\) fixes entries one and
four and acts by adjacent Hurwitz moves on the tuple
\begin{equation}\label{eq:framed}
                  (M_2,M_3,CM_5C^{-1},\ldots,CM_bC^{-1}),
                  \qquad C=M_4.
\end{equation}
This follows for the last generators because simultaneous conjugation
commutes with \(T\), and then holds for inverses as well. The product of
the entries in \eqref{eq:framed} is invariant, since \(T(a,c)\) has
product \(ac\).

We use permutations composed from right to left. The top labels on
\(2b\) symbols are
\[
\begin{aligned}
 x_0&=(0\,1)(2\,3), &x_1&=(1\,4)(2\,3),\\
 x_2&=(0\,5)(1\,6), &x_3&=(1\,2)(6\,7),\\
 x_4&=(0\,8)(7\,9), &
 x_j&=(2j-2\,\,2j)(2j-1\,\,2j+1)\quad(5\le j<b).
\end{aligned}
\]
The first four generate \(A_8\). Indeed, \(x_0x_1\) is a
three-cycle on \(\{0,1,4\}\). By conjugating that cycle by \(x_2\)
and \(x_3\), then conjugating the resulting cycles by \(x_3\) and
\(x_0\), respectively, we obtain the additional supports
\[
       \{4,5,6\},\quad \{0,2,4\},\quad \{4,5,7\},\quad \{1,3,4\}.
\]
These cycles generate \(A_8\). To see the enlargement at each step,
an alternating group on an old set of at least three symbols together
with a three-cycle on two old symbols and one new symbol generates
the alternating group on the union. Conjugations and inverses give
both orientations for every pair of old symbols, including when the
old set has size three. If the cycle has just one old symbol,
conjugate it by an old permutation moving that symbol. The product
of the original cycle with the inverse conjugate gives a cycle on
two old symbols and one new symbol, after which the other new
symbol can also be added. The displayed supports can therefore be
adjoined in order. Each \(x_j\) for \(j\ge4\) consists of two
transpositions \((a\,u)(c\,v)\) with distinct old symbols \(a,c\)
and distinct new symbols \(u,v\). Conjugate it by an even old
permutation fixing \(c\) and moving \(a\), and multiply the
two double transpositions. This gives a three-cycle on two old
symbols and \(u\). Repeat with \(a,u\) fixed and \(c\) moved to
adjoin \(v\). The full tuple generates \(A_{2b}\).

To check the coloring, let \(X=x_1,\ Y=x_2,\ C=x_3\).
Applying the crossing rule through \(U\) gives
\[
       (x_0,a,h,C,x_4,\ldots,x_{b-1}),\qquad
       a=(0\,5)(1\,4),\quad h=(1\,6)(2\,3).
\]
This four-strand calculation can also be read from the crossing
relations in \eqref{eq:prefix}. With \(z_i=Cx_{i+3}C^{-1}\) for
\(1\le i\le d-2\), the product of the middle tuple
\((a,h,z_1,\ldots,z_{d-2})\) is
\[
 p=(0\,8\,10\,\cdots\,2b-2\,\,5)
   (1\,6\,9\,11\,\cdots\,2b-1\,\,4)(2\,3).
\]
In each displayed progression the terms run only up to the indicated
endpoint. For \(b=5\) the product is the one computed below.
For larger \(b\) each additional factor inserts its two new
symbols just before \(5\) and \(4\) in the respective cycles.
The two long cycles have lengths \(d\) and \(d+1\), so
\(p\) has order \(d(d+1)\). Reading the cycles gives
\[
        pXp^{-1}=h,\qquad
        p^iYp^{-i}=z_i\ (1\le i\le d-2),\qquad p^dYp^{-d}=a.
\]
The adjacent moves in \(\delta_b\) act on \eqref{eq:framed} as the
left shift that sends the first entry \(m\) to \(p^{-1}mp\) at the
end. After \(d\) shifts all labels have been conjugated by \(p^{-1}\).
After \(d+1\) shifts the three displayed identities return the initial
middle tuple \((X,Y,z_1,\ldots,z_{d-2})\). Adding
\(d^2(d+1)k\) shifts has no effect because \(p\) has order \(d(d+1)\).

\begin{example}
At \(b=5\) the middle tuple has length three and
\(p=(0\,8\,5)(1\,6\,9\,4)(2\,3)\). Four shifts close the
coloring of \(U(\sigma_2g)^4\), and a further \(36\) shifts act
trivially. This gives the exponent \(4+36k\) on five strands.
\end{example}

The strand permutation of \(\delta_b\) is a \(d\)-cycle on positions
\(\{2,3,5,\ldots,b\}\); its exponent is \(1\bmod d\).
The permutation of \(U\) is \((1\,2)(3\,4)\). The two transpositions
insert the two omitted positions into that cycle, giving one knot
component. Its group surjects onto \(A_{2b}\) by the coloring just
constructed. A double transposition moves two independent directions in the natural
permutation representation over \(\mathbb C\). The following lower
bound uses its two transposition edges. A set of at most \(b-1\) double transpositions has a
support graph with at most \(2b-2\) edges on \(2b\) vertices. Each
generator preserves the components of the graph, so such a set
cannot generate a transitive group. The quotient forces
\(\mr(K_{b,k})\ge b\), since every positive meridian is conjugate to
a top label. A closed \(b\)-strand braid gives a bridge presentation
with at most \(b\) maxima, and the Wirtinger relations of a bridge
presentation give one generating meridian per bridge. Thus
\(\mr(K_{b,k})=\br(K_{b,k})=b\).
